\chapter{Exercises}
\section{List 1}
\begin{dang}
    \textbf{Problem 1}. State the definition of monomorphism, 
    and prove that an injective function of sets is a monomorphism in Set.
\end{dang}
Solution: A monomorphism $f\in\text{Hom}_{\callC}(B,C)$ is a morphism that is left cancelative in the sense that 
for any object $A$ and morphisms $g,h\in\text{Hom}_{\callC}(A,B)$, we have:
\begin{equation}
    fg = fh \Rightarrow g= h.
\end{equation}
Now, if $\callC = \textbf{Set}$, the equality $fg=fh$ implies $f(g(a))=f(h(a)),\;\forall a\in A$ and if $f$ injective, by 
definition, we have $g(a)=h(a),\;\forall a\in A$, i.e., $g=h$.

\begin{dang}
    \textbf{Problem 2}. Prove that $h:\mathbb{Z}\hookrightarrow \mathbb{Q}$ is an epimorphism in Ring.
\end{dang}
Solution: Let $f,g:\bbQ\to R$ be two ring morphisms such that $fh=gh$, i.e., $f(n)=g(n),\;\forall n\in\bbZ$. Here, we 
assume that a morphism in Ring send unity to unity, so $f(1)=g(1)=1_R$. Furthermore, since $\bbQ$ is a field, 
$f:\bbQ\hookrightarrow R$ and $f(\bbQ^\times)\subset R^\times$. In particular $m1_R\in R^\times$. Now we have
\begin{equation}
    m1_R(f(\frac{1}{m})-g(\frac{1}{m})) = 0\Rightarrow f(\frac{1}{m})=g(\frac{1}{m}).
\end{equation}
This proves that $f=g$.


\begin{dang}
    \textbf{Problem 3}. Write the proof that any two terminal objects in a category 
    $\mathcal{C}$ are isomorphic to each other.
\end{dang}
Solution: Let $T,T^\prime$ be two terminal objects in $\mathcal{C}$. Let $f\in\Hom_{\callC}(T,T^\prime)$ 
and $T^\prime\in\Hom_{\callC}(T^\prime,T)$, then $f\circ g\in\Hom_{\callC}(T,T)=\{\mathds{1}_{T}\}$. Therefore 
$fg = \mathds{1}_{T}$ and in a similar way we see that $gf = \mathds{1}_{T^\prime}$. This proves that 
$f$ is an isomorphism.

\begin{dang}
    \textbf{Problem 4}. State an "explicit" definition of the free product \(G*H\) of groups \(G\) and \(H\). 
Prove that the free product is a coproduct in Grp.
\end{dang}
Solution: Let $A = G \textstyle\coprod H$, we define
\begin{equation}
    G*H = \coprod_{n\geq1}A^n\bigg/\sim 
\end{equation}
where we the equivalence relation is defined by 
\begin{equation}
    \begin{aligned}
        &(a_1,\dots, a_i,a_{i+1},\dots,a_n)\sim (a_1,\dots, a_ia_{i+1},\dots,a_n)
        \text{ if }a_1,a_2\in G\text{ or } a_1,a_2\in H\\
        &(a_1,\dots, a_{i-1},a_i,a_{i+1} \dots, a_n)\sim (a_1,\dots, a_{i-1},a_{i+1} \dots, a_n)\;\text{ if } 
        a_i = 1_G \text{ or } a_i = 1_H
    \end{aligned}
\end{equation}
The multiplication is defined by 
\begin{equation}
    [(a_1,\dots,a_n)]* [(b_1,\dots, b_m)] = [(a_1,\dots,a_n,b_1,\dots, b_m)].
\end{equation}
By the equivalence relations defined above, the multiplication rule is clearly well defined.\par  
We see that $(G*H, *)$ is a group. The associativity and existence of unity is clear and 
the inverse operation is defined by 
\begin{equation}
    [(a_1,a_2,\dots,a_n)]^{-1}=[(a_n^{-1},\dots, a_2^{-1},a_1^{-1})].
\end{equation}
Now we define morphims
\begin{equation}
    i_1: G \to G*H: g\mapsto [(g)],\; i_2:H\to G*H: h\mapsto [(h)].
\end{equation}
We claim that $(G*H, i_1,i_2)$ is a coproduct of $G$ and $H$. Suppose that we have morphisms 
$f_1: G\to K, f_2:H\to K$, then we define a morphism 
\begin{equation}
  f:  G*H \to K: [(a_1,\dots, a_n)]\mapsto  f_{j_1}(a_1)\cdots f_{j_n}(a_n)
\end{equation}
where we let $j_k=1\;(\text{resp. }2)$ if $a_k\in G\;(\text{resp. }H)$. Furthermore, we see that 
$f i_j=f_j,\;j=1,2$. For uniqueness, we observe that if $h:G*H\to K$ is a morphism, such that 
$h i_j = f_j,\;j=1,2$, then $h([(g)])=f_1(g)$ and $h([(h)])=f_2(h)$. So 
\begin{equation}
    h([(a_1,\dots, a_n)])=h([(a_1)]*\cdots*[(a_n)])=h([a_1])\cdots h([(a_n)])
\end{equation}
and we see that $h([(a_k)])=f_{j_k}(a_k)$ as defined above, therefore $h=f$.



\begin{dang}
\textbf{Problem 5}. Let $\mathcal{C}$ be a pre-additive category and $X,Y\in \mathcal{C}$ . 
Show that if a product $X\times Y$ of $X$ and $Y$ exists, then $X\times Y$ is also a coproduct of $X$ and $Y$ .   
\end{dang}
Solution: A product $X\times Y$ comes with projections $p_1:X\times Y\to X$ and $p_2:X\times Y\to Y$ such that the following universal
property is satisfied. 
\begin{center}
    \begin{tikzcd}
         &Z\arrow[d,dashed,"\exists !h"]\arrow[ddl,bend right = 30, swap,"f_1"]\arrow[ddr,bend left = 30, "f_2"]&\\
         &X\times Y\arrow[dl,swap,"p_1"]\arrow[dr,"p_2"]&\\
         X&&Y
    \end{tikzcd}
\end{center}
So we define $i_1:X\to X\times Y$ and $i_2:Y\to X\times Y$ using this universal property:
\begin{center}
    \begin{tikzcd}
        &X\arrow[d,dashed,"\exists !i_1"]\arrow[ddl,bend right = 30, swap,"\mathds{1}_X"]
        \arrow[ddr,bend left = 30, "0_{XY}"]&\\
         &X\times Y\arrow[dl,swap,"p_1"]\arrow[dr,"p_2"]&\\
         X&&Y
    \end{tikzcd}

    \begin{tikzcd}
        &Y\arrow[d,dashed,"\exists !i_1"]\arrow[ddl,bend right = 30, swap,"0_{YX}"]
        \arrow[ddr,bend left = 30, "\mathds{1}_Y"]&\\
         &X\times Y\arrow[dl,swap,"p_1"]\arrow[dr,"p_2"]&\\
         X&&Y
    \end{tikzcd}
\end{center}
So we have the following properties 
\begin{equation}\label{equation: relations between projections and injections}
    \left\{\begin{aligned}
        &\mathds{1}_X=p_1i_1,\; p_2i_1 = 0_{XY}\\
        &\mathds{1}_Y = p_2i_2,\;p_1i_2 = 0_{YX}
    \end{aligned}\right.
\end{equation}
Now, we claim that 
\begin{equation}
    i_1p_1+i_2p_2 = \mathds{1}_{X\times Y}\quad(*)
\end{equation}
which can be seen by the diagram below and the universal 
property of the product.
\begin{center}
    \begin{tikzcd}
        &X\times Y\arrow[d,dashed,"i_1p_1+i_2p_2"]\arrow[ddl,bend right = 50, swap,"p_1"]
        \arrow[ddr,bend left = 50, "p_2"]&\\
         &X\times Y\arrow[dl,swap,"p_1"]\arrow[dr,"p_2"]&\\
         X&&Y
    \end{tikzcd}
\end{center}
Now, suppose that we have morphisms $f_1:X\to Z, f_2:Y\to Z$. Then if we define 
$h=f_1p_1+f_2p_2$, the diagram below commutes.
\begin{center}
    \begin{tikzcd}
        &Z&\\  
        &X\times Y\arrow[u, dashed, "h"]&\\ 
        X\arrow[ur, swap,"i_1"]\arrow[uur, bend left = 30, "f_1"]&&Y\arrow[ul, "i_2"]\arrow[uul,swap, bend right = 30, "f_2"]
    \end{tikzcd}
\end{center}
Furthermore, relations~(\ref{equation: relations between projections and injections}) and equation (*) guarantees that 
this is the unique expression possible for $h$ that makes the above diagram commutes.


\begin{dang}
    \textbf{Problem 6}. Show that existence of binary products implies existence of all finite \(n\)-ary products 
    for \(n\geq 2\).
\end{dang}
Solution: By a finite $n$-ary product $(n\geq2)$, we mean an object $X=\prod_{i=1}^n X_i$ with morphims 
$p_i: X\to X_i$ satisfying the following universal property. 
Take a colection of morphims $f_i:Y\to X_i,\;(i=1,\dots,n)$, then there exist an unique morphism $f: Y\to X$ such that 
$p_if = f_i,\;\forall i$. By induction, assume that there exist all finite $(n-1)-$ary products and let 
$X^\prime =\prod_{i=1}^{n-1}X_i$. We want to show that $X=X^\prime \times X_n$ satisfies the universal property for 
a $n-$ary product. The projections are obtained as follows: we have the projection 
$p_n:X\to X_n$ and for $i<n$ we compose the projection 
$q:X\to X^\prime$ with the $i$th projection $p_i^\prime:X^\prime\to X_i$ to get $p_i:X\to X_i$.\par 

Now assume we have morphisms $f_i:Y\to X,\;(i=1,\dots,n)$. Then, by the universal property of the finite 
$(n-1)-$ary product, there exists a (unique) morphism $f^\prime:Y\to X^\prime$ such that 
$p_i^\prime f^\prime = f_i,\; i<n$. Now we have the diagram:
\begin{center}
    \begin{tikzcd}
        &Y\arrow[d, dashed, "\exists!h"]\arrow[ldd,swap, bend right = 30, "f^\prime"]\arrow[rdd, bend left = 30, "f_n"]&\\
        &X\arrow[ld,"q"]\arrow[dr,swap, "p_n"]&\\
        X^\prime&& X_n
    \end{tikzcd}
\end{center}
Now we observe that 
\begin{equation}
    p_ih=(p_i^\prime q)h=p_i^\prime(qh)=p_i^\prime f^\prime = f_i,\; i<n.
\end{equation}
This show existence of $h$. Now, for uniqueness, suppose that $h^\prime$ is such that 
$p_i h^\prime=f_i,\;\forall i$. Then, for $i<n$ we have $p_i^\prime(h^\prime q)=f_i$. So by the universal property of the 
finite $(n-1)-$ary product $X^\prime$, we must have $h^\prime q = f^\prime$. Then, by applying the universal 
property of the product $X=X^\prime\times X_n$, we must have $h^\prime = h$.

\begin{dang}
    \textbf{Problem 7}. Consider a composition
\[
A\longrightarrow B\longrightarrow C
\]
in an abelian category, in which \(g\circ f = 0\). Using only universal properties, 
show that there is a canonical morphism from \(\operatorname{im}(f)\) to \(\ker(g)\). 
[Recall \(\operatorname{im}(f) = \ker(\operatorname{coker}(f))\) by definition.]
\end{dang}
Solution: From the universal property of the cokernel we have the diagram.
\begin{center}
    \begin{tikzcd}
        A\arrow[rr,"f"]\arrow[dd]&&B\arrow[dd,"g"]\arrow[ld,"p"]\\
        &\text{Coker}(f)\arrow[rd,dashed,"\exists!\tilde{g}"]&\\
        0\arrow[ur]\arrow[rr]&&C
    \end{tikzcd}
\end{center}
Now, by the universal property of the kernel, we have:
\begin{center}
    \begin{tikzcd}
        A\arrow[rr,"f"]\arrow[dd]\arrow[dr,dashed,swap,"\exists!\tilde{f}"]&&B\arrow[dd,"p"]\arrow[dr,"g"]&\\
        &\text{Ker}(p)=\text{Im}(f)\arrow[dl]\arrow[ru,swap,"i"]&&C\\
        0\arrow[rr]&&\text{coker}(f)\arrow[ru,swap,"\tilde{g}"]&
    \end{tikzcd}
\end{center}
Finally, by applying the universal property of the kernel again, we have:
\begin{center}
    \begin{tikzcd}
        \text{Im}(f)\arrow[dd]\arrow[rr,"i"]\arrow[rd,dashed,"\exists!\gamma"]&&B\arrow[dd,"g"]\\
        &\text{Ker}(g)\arrow[ru]\arrow[ld]&\\
        0\arrow[rr]&&C
    \end{tikzcd}
\end{center}
as desired.

\begin{dang}
    \textbf{Problem 8}. Let \(G:\mathbf{Ab}\to \mathbf{Set}\) be the forgetful functor, and let \(F:\mathbf{Set}\to \mathbf{Ab}\) 
    be the functor of forming a free abelian group. Show that \(F\) is left adjoint to \(G\).
\end{dang}
Solution: We have the map 
\begin{equation}
    \text{Hom}_{\textbf{Ab}}(F(X),Y) \rightarrow \text{Hom}_{\textbf{Set}}(X,G(Y)): f\mapsto f|_X
\end{equation}
Here, we let 
$F(X) = \{f:X\to \bbZ: f(x)\neq 0\text{ for a finite number of } x\in X\}$ and denote
$\hat{x}\in F(X)$ for the map defined by $\hat{x}(y)=\delta_{xy}$. So we see that the 
set theoretic map $\iota:X\to F(X):x\to \hat{x}$ is an injection and this justifies 
the notation $f|_X$ used in $(*)$. Moreover, show that this map is an isomorphism is 
equivalent to show that we have the following universal property: For any set map of sets 
$f:X\to G(Y)$, there exists a unique morphism $\hat{f}:F(X)\to Y$ such that $\hat{f}\iota = f$. So we have no choice but 
define $\hat{f}(\hat{x})=f(x)$ and extend to $F(X)$ by $\bbZ-$linearity, this defines $\hat{f}$ uniquely.



\begin{dang}
\textbf{Problem 9}. Compute $(\mathbb{Z} / 5)\otimes_{\mathbb{Z}}(\mathbb{Z} / 6)$ 
and $(\mathbb{Z} / p^{2})\otimes_{\mathbb{Z}}(\mathbb{Z} / p)$ where $p$ is a prime.
\end{dang}
Solution: First we claim that $(\mathbb{Z} / 5)\otimes_{\mathbb{Z}}(\mathbb{Z} / 6)=(0)$. In fact, if 
$a\otimes b\in (\mathbb{Z} / 5)\otimes_{\mathbb{Z}}(\mathbb{Z} / 6)$, then since $6^2\equiv1\mod5$, we have
\begin{equation}
    a\otimes b = a\cdot 6^2\otimes b = a\cdot6\otimes 6\cdot b = a\cdot6\otimes 0=0.
\end{equation}
For the second group, we claim that $(\mathbb{Z} / p^{2})\otimes_{\mathbb{Z}}(\mathbb{Z} / p)\cong\bbZ/p$. To show this, 
we show that $\bbZ/p$ satisfies the universal property of the tensor product (here we will use the notation of 
problem 10). Let $\varphi_0:\bbZ/p^2\times\bbZ/p\to \bbZ/p:([m],[n])\mapsto [mn]$, it's a well defined balanced map, i.e., 
$\varphi_0\in\text{Bal}_{A,B}(\bbZ/p)$. In particular $\bbZ/p$ is generated by the image of $\varphi_0$. Now, let 
$f\in \text{Bal}_{A,B}(X)$, we want to define $\tilde{f}$ in such a way the following diagram commutes.
\begin{center}
    \begin{tikzcd}
    \bbZ/p^2\times\bbZ/p\arrow[d,swap,"f"]\arrow[r,"\varphi_0"]&\bbZ/p\arrow[ld,"\tilde{f}"]\\
    X&
    \end{tikzcd}
\end{center}
So we set $\tilde{f}([n])=f(1,[n])$, which is well defined while makes the diagram commutes.

\begin{dang}
\textbf{Problem 10}. Let $R$ be a unital ring. In order to formulate the definition of tensor product of $R$-modules, 
we introduce a definition:\par 
\textbf{definition 0.1}. Let $A$ be a right $R$-module and $B$ a left $R$-module, and $X$ an abelian group. 
A balanced map $f\in \operatorname{Bal}_{A,B}(X)$ is a function $f:A\times B\to X$ such that

\[
f(a_{1} + a_{2},b) = f(a_{1},b) + f(a_{2},b)\quad \mathrm{for~all~}a_{1},a_{2}\in A\mathrm{~and~}b\in B
\]
\[
f(a,b_{1} + b_{2}) = f(a,b_{1}) + f(a,b_{2})\quad \mathrm{for~all~}a\in A\mathrm{~and~}b_{1},b_{2}\in B
\]
\[
f(ar,b) = f(a,rb)\quad \mathrm{for~all~}a\in A,b\in B\mathrm{~and~}r\in R.
\]

Now a tensor product of a right $R$-module $A$ with a left $R$-module $B$ is an abelian group $A \otimes_R B$ 
together with a $\varphi_0 \in \operatorname{Bal}_{A,B}(A \otimes_R B)$, such that for any balanced map $f$ as below, 
there exists a unique linear map $\tilde{f}$ making the diagram commute.

\begin{center}
    \begin{tikzcd}
    A\times B\arrow[dr,swap,"f"]\arrow[rr,"\varphi_0"] & & A\otimes_{R}B\arrow[ld,"\tilde{f}"]\\ 
    & X&
    \end{tikzcd}
\end{center}

Verify that

\[
A\otimes_{R}B\cong Q = \frac{A\otimes_{\mathbb{Z}}B}{\langle ar\otimes b - a\otimes rb\mid a\in A,b\in B,r\in R\rangle}.
\]
\end{dang}
Solution:  To prove the asserted isomorphism, we show that the quotient of abelian groups $Q$ satisfies the universal 
property of the tensor product $A\otimes_RB$. Let $f\in\text{Bal}_{A,B}(X)$, then $f$ is balanced by viewing $A,B$ 
as abelian groups. By the universal property of $A\otimes_{\bbZ}B$, there exist maps
$\hat{\varphi}_0$ and $\hat{f}$, such that:
\begin{center}
    \begin{tikzcd}
    A\times B\arrow[dr,swap,"f"]\arrow[rr,"\hat{\varphi}_0"] & & A\otimes_{\bbZ}B\arrow[ld,"\hat{f}"]\\ 
    & X&
    \end{tikzcd}
\end{center}
Now, since $f\in \text{Bal}_{A,B}(X)$, we have 
\begin{equation}
    \Ker(\hat{f})\supset \langle ar\otimes b - a\otimes rb\mid a\in A,b\in B,r\in R\rangle.
\end{equation}
Then, there exists $\tilde{f}$ such that the following diagram commutes:\begin{center}
    \begin{tikzcd}
    A\times B\arrow[dr,swap,"f"]\arrow[rr,"\hat{\varphi}_0"] & & A\otimes_{\bbZ}B\arrow[ld,"\hat{f}"]\arrow[r,"p"]&
    Q \arrow[lld,bend left = 20,"\tilde{f}"]\\ 
    & X& &
    \end{tikzcd}
\end{center}
Now, by construction, the quotient $Q$ is generated by the image of the map 
$\varphi_0 = p\hat{\varphi}_0$ and since $f = \tilde{f}\varphi_0$ we see that $\tilde{f}$ is uniquely determined 
by $f$. Therefore $(Q,\varphi_0)$ satisfies the universal property for the tensor product $A\otimes_RB$ and 
$Q\cong A\otimes_RB$.






\section{List 2}
\begin{dang}
\textbf{Problem 1}.
Let $A$ be an abelian group, and consider the action
\[
A \times A \longrightarrow A
\]
given by addition in the group. Does this action make $A$ into an
$A$-module? Justify your answer.
\end{dang}
Solution: For $a,m\in A$ we define $am = a+m$. One of the module axioms says 
$(a+a^\prime)m = am + a^\prime m$, so we must have 
\begin{equation}
    a+a^\prime + m = a+m + a^\prime + m\Rightarrow m=0.
\end{equation}
So, the given action defines an $A-$module if and only if $A$ is the trivial group.\qed 


\begin{dang}
    \textbf{Problem 2}.
The homologies $H_n(G,A)$ of a group $G$ with coefficients in the
$G$-module $A$ are the derived functors of the coinvariants functor
\[
(-)_G = \mathbb{Z} \otimes_{\mathbb{Z}G} (-).
\]
(where $\mathbb{Z}$ is the trivial $G$-module). Using a convenient
projective resolution of $\mathbb{Z}$, or otherwise, compute group
homology $H_n(G,A)$ where $G$ is the cyclic group $\mathbb{Z}/n$ and
$A$ is a trivial $G$-module.
\end{dang}
Solution: We use the periodic projective resolution 
\begin{equation}
    \cdots\overset{g-1}{\to} \bbZ G\overset{N}{\to}\bbZ G\overset{g-1}{\to}\bbZ G\overset{\epsilon}{\to}\bbZ 
\end{equation}
where $G=C_n=\langle g:g^n=1\rangle$, $\epsilon:g\mapsto 1$ and $N=1+g+\cdots+g^{n-1}$. Now we apply  
$(-)\otimes_{\bbZ G}A$ (which is left exact) to the above sequence to get 
\begin{equation}
    \cdots\overset{g-1}{\to} A\overset{N}{\to}A\overset{g-1}{\to}A\overset{\epsilon}{\to}A_G.
\end{equation}
Here we used the general fact that if $M$ is a $R$-module, then $R\otimes_R M\cong M$ as $R-$modules. 
Now we have
\begin{equation}
    \begin{split}
        H_n(G,A) &= [R^n\bbZ\otimes_{\bbZ G}(-)](A)\quad(\text{definition})\\ 
        &=[L_n(-)\otimes_{\bbZ}A](\bbZ)\quad(\text{balancing})\\
        &= H_n(\cdots\overset{g-1}{\to} A\overset{N}{\to}A\overset{g-1}{\to}A\to0).
    \end{split}
\end{equation}
Now, if $A$ is a trivial module, then $ga=a,\forall g\in G$. Therefore $\text{Ker}(g-1)=A$ and 
$\text{Im}(N)=nA$. So we get 
\begin{equation}
   H_m(G,A) =\left\{\begin{aligned}
        &A/nA,\;m \text{ odd};\\
        &\{a\in A:na = 0\},\;m\text{ even }>0;\\
        &A_G,\;m=0
    \end{aligned}\right.
\end{equation}
\qed



\begin{dang}
    \textbf{Problem 3}.
    \textbf{The Curry adjunction}:
Let $A \in \mathbf{Set}$ be a fixed set, and let
$L : \mathbf{Set} \to \mathbf{Set}$ and
$R : \mathbf{Set} \to \mathbf{Set}$ be the functors
\[
L(X) = X \times A
\qquad\text{and}\qquad
R(X) = \operatorname{Hom}_{\mathbf{Set}}(A,X).
\]
Prove that $L$ is left adjoint to $R$, i.e., construct a natural
isomorphism
\[
\operatorname{curry} :
\operatorname{Hom}_{\mathbf{Set}}(X \times A,Y)
\longrightarrow
\operatorname{Hom}_{\mathbf{Set}}
\bigl(X,\operatorname{Hom}_{\mathbf{Set}}(A,Y)\bigr).
\]
\end{dang}
Solution: Let $X\in\catset$, we claim that $(X\times A, u_X),\;u_X:X\to R(X\times A):x\mapsto [a\mapsto (x,a)]$ is a universal 
from $X$ to $R$. Let $f\in\Hom_{\catset}(X,R(Y))$, if we define $\tilde{f}:X\times A\to Y:(x,a)\mapsto f(x)(a)$, then 
the diagram below commutes 
\begin{center}
    \begin{tikzcd}
      X\arrow[r,"u_X"]\arrow[d,"f"] &  R(X\times A)\arrow[dl,dashed,"R(\tilde{f})"]    \\   R(Y)&
    \end{tikzcd}
\end{center}
In fact $R(\tilde{f})u_X=\Hom(A,\tilde{f})u_X=\tilde{f}\circ u_X$ and 
$\tilde{f}\circ u_X(x)(a)=\tilde{f}(x,a)=f(x)(a)$. In particular, we see that commutativity of the diagram 
uniquely defines $\tilde{f}$. Then we see that for each $X\in\catset$, there exist 
a universal $(LX, u_X)$ from $X$ to $R$. This statement is equivalent to say that the 
map $\eta_{X,Y}:\Hom_{\catset}(LX,Y)\rightarrow \Hom_{\catset}(X,RY): g\mapsto R(g)u_X$ is an isomorphism, i.e., 
the assigment $Y\mapsto \eta_{X,Y}$ gives a natural isomorphism 
$\Hom_{\catset}(LX,-)\Rightarrow \Hom_{\catset}(X,R-)$. This proves that $L$ is left adjoint to $R$.\qed



\begin{dang}
    \textbf{Problem 4}.
Prove the following lemma from class.

\medskip

\noindent
\textbf{Lemma 0.1.}
If $L : \mathcal{A} \to \mathcal{B}$ is a functor with right adjoint
$R$, and
\[
\eta : \operatorname{Id} \Rightarrow R \circ L
\qquad\text{and}\qquad
\varepsilon : L \circ R \Rightarrow \operatorname{Id}
\]
are as in class (see Corollary 3.4.9). Then
\[
\varepsilon_{L(X)} \circ L(\eta_X) = \operatorname{Id}_{L(X)}.
\]

\medskip

\noindent
\textit{Hint:}
There are several equivalent definitions of adjunction. The one I have
given is in terms of a natural isomorphism
\[
\beta_{X,Y} :
\operatorname{Hom}(L(X),Y)
\longrightarrow
\operatorname{Hom}(X,R(Y)).
\]
There are other definitions, in particular one that is in terms of two
natural transformations
\[
\eta : \operatorname{Id}_{\mathcal A} \Rightarrow R \circ L,
\qquad
\varepsilon : L \circ R \Rightarrow \operatorname{Id}_{\mathcal B}.
\]

One of the axioms in this second definition is
\[
\begin{tikzcd}[column sep=large, row sep=large]
Y
    \arrow[r, "L(\eta)"]
    \arrow[dr, "\operatorname{Id}_Y"']
&
LRY
    \arrow[d, "\varepsilon_Y"]
\\
&
Y
\end{tikzcd}
\]

The identity to be proved follows fairly easily from this axiom. So the
exercise, really, is to understand this second definition of adjoint
functors, and why it follows from the first.
\end{dang}
Solution: In the context of the Yoneda's lemma, we have a bijection 
\begin{equation}
  A\to  \{\eta: \Hom_{\callC}(A,-)\Rightarrow F\}:a\mapsto [a_B: \Hom_{\callC}(A,B)\to F(B):k\mapsto F(k)a].
\end{equation}
Its inverse is given by $\eta_A\mapsto \eta_A(\mathds{1}_A)$. In this context we say that 
$(A,a),\;a=\eta_A(\mathds{1}_A)$ represents $\eta_A$. Now, coming back to the problem, we have a 
natural isomorphism $\eta_{LX}:\Hom_{\callA}(LX,-)\Rightarrow\Hom_{\callB}(X,R-)$ given by the assigment 
$Y\mapsto \eta_{X,Y}$. So, by the Yoneda's lemma, we see that the pair 
$(LX, u_X),u_X=\beta_{X,LX}(\mathds{1}_{LX})$ represents $\eta_{LX}$. So we see that $\beta_{X,Y}$ is given by the formula 
\begin{equation}
    k\mapsto \Hom(X,R(k))(u_X)= R(k)u_X.
\end{equation}
The fact that this map is a bijection, implies that $(LX, u_X)$ satisfies the following statement: Let 
$g\in\Hom_{\callB}(X,RY)$, then there exists a unique morphism $\tilde{g}\in\Hom_{\callA}(LX,Y)$ such that the diagram below
commutes. 
\begin{center}
    \begin{tikzcd}
      X\arrow[r,"u_X"]\arrow[d,swap,"g"]  &  RLX\arrow[dl,dashed,"R(\tilde{g})"]   \\  RY &
    \end{tikzcd}
\end{center}
So we conclude that $(LX,u_X)$ is a universal from $X$ to $R$. Now, for every $Y\in\callA$, define 
$v_Y: LRY\to Y$ as the unique morphism making the diagram below commutes.
\begin{center}
    \begin{tikzcd}
        RY\arrow[r,"u_{RY}"]\arrow[d,swap,"\mathds{1}_{RY}"]  &  RLRY\arrow[dl,dashed,"R(v_Y)"]  \\   RY &
    \end{tikzcd}
\end{center}
We claim that the pair $(RY, v_Y)$ is a universal from $L$ to $Y$, i.e., for every $f\in\Hom_{\callA}(LX, Y)$ 
there exists a unique morphism $\tilde{f}\in\Hom_{\callB}(X,RY)$ such that 
\begin{center}
    \begin{tikzcd}
       Y & LRY\arrow[l,swap,"v_Y"] \\ LX\arrow[u,"f"]\arrow[ru,dashed,swap,"L(\tilde{f})"]&
    \end{tikzcd}
\end{center}
This is equivalent to prove that $\xi_{X,Y}:\Hom_{\callB}(Y,RX)\to\Hom_{\callA}(LY,X):f\mapsto v_YL(f)$ is 
a bijection. In fact we have 
\begin{equation}
    \begin{split}
        \beta_{X,Y}\xi_{X,Y}(f)&=R(v_YL(f))u_X\\
        &=R(v_Y)RL(f)u_X\\
        &=R(v_Y)u_{RX}f\\
        &=\mathds{1}_{RY}f\\
        &=f.
    \end{split}
\end{equation}
Then we see that $\xi_{X,Y}=\beta_{X,Y}^{-1}$.\qed 

\begin{dang}
    \textbf{Problem 5}.
Deduce the classical version of Hilbert 90 from the more general
cohomological statement, using the ``convenient complex'' for
$H^\bullet(\mathbb{Z}/n,-)$ that we met in class:

\medskip

\noindent
\textbf{Theorem 0.2.}
Let $L/K$ be a Galois extension, and suppose that
\[
G = \operatorname{Gal}(L/K)
\]
is cyclic,
\[
G = \langle \sigma \rangle \cong \mathbb{Z}/n.
\]
Then for every $\alpha \in L$ of norm $N(\alpha)=1$, there exists
$\beta \in L$ such that
\[
\alpha = \frac{\sigma(\beta)}{\beta}.
\]
\end{dang}
Solution: The cohomological statement of Hilbert 90's is $H^1(G,L^\times)=0$, where $G=\text{Gal}(L|K)$ is the Galois 
group of the Galois extension $L|K$. Take the periodic projective resolution of $\bbZ$.
\begin{equation}
    \bbZ G\overset{N}{\to}\bbZ G \overset{\sigma-1}{\to}\bbZ G \overset{\epsilon}{\to}\bbZ\to0.
\end{equation}
Now, apply the left exact contravariant functor $\Hom(-,L^\times)$ to the cochain complex:
\begin{equation}
    0\to \Hom_{\bbZ G}(\bbZ G,L^\times)\overset{-\circ(\sigma-1)}{\to} \Hom_{\bbZ G}(\bbZ G,L^\times)
    \overset{-\circ N}{\to}\cdots 
\end{equation}
Now we have an isomorphism $\Hom_{\bbZ G}(\bbZ G,L^\times)\to L^\times: f\mapsto f(1)$, with inverse given by 
$a\mapsto [f_a:g\mapsto g(a)]$. Moreover, we see that 
$f_a\circ (\sigma-1)(1)=f_a(\sigma-1)=(\sigma - 1)f_a(1)=(\sigma-1)a=\sigma(a)a^{-1}$ and 
$f_a\circ N(1) = f_a(N)=Nf_a(1)=a\sigma(a)\cdots\sigma^{n-1}(a)=N_{L|K}(a)$. So, the statement $H^1(G,L^\times)=0$, means that 
$\text{Im}(\sigma-1)=\text{Ker}(N_{L|K})$. Therefore, if $N_{K|L}(a)=1$, then there exists $\beta\in L^\times$ 
such that $\alpha = (\sigma-1)\beta=\sigma(\beta)\beta^{-1}$.\qed
\begin{dang}
    \textbf{Problem 6}.
Let $G$ be a group and $N \subset G$ a normal subgroup. If $A$ is a
$G$-module, verify that the invariants $A^N$ carries a natural action
of the quotient group $G/N$. Thus we may consider the functor
\[
(-)^N :
G\text{-mod}
\longrightarrow
(G/N)\text{-mod}.
\]
\end{dang}
Solution: The action $gN\cdot a=ga$, defines $G/N-$module structure on $A^N$. It's well defined because 
$gN=g^\prime N\Rightarrow g^{-1}g^\prime = n\in N$. Therefore $g^{-1}g^\prime a = a\Rightarrow g^\prime a=ga$. Furthermore 
$ga\in A^N,\;\forall g\in G$ because if $n\in N$, then there exists 
$n^\prime\in N$ such that $ng = gn^\prime$, so $nga = gn^\prime a = ga$.\qed

\begin{dang}
    \textbf{Problem 7}.
Let $F : \mathcal{A} \to \mathcal{B}$ be a left exact functor, and
$E : \mathcal{B} \to \mathcal{C}$ an exact functor. Assume
$\mathcal{A}$ has enough injectives. Prove that
\[
\bigl[R^n(E \circ F)\bigr](A)
\cong
\bigl[E \circ R^nF\bigr](A)
\]
for all $A \in \mathcal{A}$. Deduce that there exists a natural action
of $G/N$ on the cohomology
\[
H^\bullet(N,A)
\]
(where $G,N,A$ as in Exercise 6).
\end{dang}
Solution: 

\begin{dang}
    \textbf{Problem 8}.
Let $G$ be a group and $N \subset G$ a normal subgroup, and $A$ a
$G$-module. Consider the bar complex, whose terms are
\[
C^p(N,A) = \operatorname{Maps}(N^p \to A).
\]
Check that
\[
[g \cdot \varphi](n_1,\ldots,n_p)
=
g\varphi
\bigl(g^{-1}n_1g,\ldots,g^{-1}n_pg\bigr),
\]
for all $p \in \mathbb{Z}_+$, defines a $G$-action, and a
$G/N$-action, on
\[
H^\bullet(N,A).
\]
\end{dang}
Solution: Let $\callD$ be the category of pairs $(G,A)$, where $A$ is a $G-$module, with morphisms given by 
$(\alpha,f):(G,A)\to (G^\prime, A^\prime)$, where $\alpha:G\to G^\prime$ and $f:A^\prime\to A$ are group 
morphisms such that
\begin{equation}
    f(\alpha(g)a)=g f(a),\;\forall g\in G, a\in A^\prime\quad(*).
\end{equation}
Now we claim that 
$H^k:\callD\to \catab: (G,A)\mapsto H^k(G,A)$ is a functor. So we need to prove that a morphism 
$(\alpha,f)$ induces a morphism $H^k(\alpha,f): H^k(G,A)\to H^k(G^\prime,A^\prime)$. 
Let $P_k=\bbZ G^{k+1}, P_k^\prime =\bbZ (G^\prime)^{k+1}$, so $P_\bullet = (P_k)_{k\geq0}$ 
and $P_\bullet^\prime = (P_k^\prime)_{k\geq0}$ are projective resolutions of $\bbZ$ viewed as a trivial module of 
$G$ and $G^\prime$ respectively. Using $\alpha$, we can regard both chains as chains of $G$-modules. However, the second 
one is no more projective as a $G$-module, but it's still acyclic. So we can shift the identity $\bbZ\to \bbZ$ 
to a chain morhism $\tau: P_\bullet\to P_\bullet^\prime$, so it satisfies 
$\tau_k(gm)=\alpha(g)t_k(m),\;m\in P_k$ (also note that any two choice of liftings are homotopic). Now we define the map 
\begin{equation}
    \phi_k:\Hom_{\bbZ G^\prime}(P_k^\prime, A^\prime)\to\Hom_{\bbZ G}(P_k,A):\varphi\mapsto f\circ\varphi\circ\tau_k.
\end{equation}
We observe that 
\begin{equation}
    \begin{split}
        \phi_k(\varphi)(gm)&=f\varphi\tau_k(gm)\\
        &=f\varphi(\alpha(g)\tau_k(m))\\
        &=f(\alpha(g)\varphi\tau_k(m))\\
        &=g\phi_k(\varphi)(m)\quad(\text{because of }(*))
    \end{split}
\end{equation}
Also, we observe that the differentials in the cochain complex $(\Hom_{\bbZ G}(P_k,A))_{k\geq0}$ are given by 
$\Hom(d,A)=(-)\circ d$, where $d(g_0,\dots, g_k)=\sum_i (-1)^i(g_0,\dots,\hat{g}_i,\dots,g_k)$. Using this expression 
it's easy to see that $d\circ\phi_k = \phi_k\circ d$. Therefore, $\phi_k$ descend to morphisms in cohomology:
\begin{equation}
    \Hom(f,\tau): H^k(G^\prime,A^\prime)\rightarrow H^k(G,A):[\varphi]\mapsto [f\varphi\tau_k]\quad(**).
\end{equation}
Now, let $g\in G$ and $N\subset G$ a normal subgroup. Then we have a morphism $(\alpha,f):(N,A)\to (N,A)$
given by $\alpha:N\to N:n\mapsto gng^{-1}$ and $f(a)=g^{-1}a$. The axiom $(*)$ is thus verified:
\begin{equation}
   f(\alpha(n)a)=f(gng^{-1}a)=ng^{-1}(a)=nf(a).    
\end{equation}
Then we have an induced $G$ action on $H^k(N,A)$. In this case, we can take the lifting $\tau$ 
defined by $\tau_k(m)=gm$. It satisfies 
\begin{equation}
       \tau_k(hm)=ghm = ghg^{-1}gm = \alpha(h)\tau_k(m),\;\forall h\in G. 
\end{equation}
Therefore the associated action on $H^k(N,A)$ is given by $c(g)^*:[\varphi]\mapsto [m\mapsto g^{-1}\varphi(gm)]$ 
because of formula $(**)$. If $g\in N$, then $c(g)^*$ is the identity map, because $\varphi$ is 
a morphism of $N-$modules. It remains to prove that under the bijection 
$C^k(N,A)\cong \Hom_{\bbZ N}(\bbZ N^{k+1},A)$ the $G-$action agrees with the one defined in the problem.

\begin{dang}
    \textbf{Problem 9}.
We continue in the context of the previous exercise. But now $G$ is
the symmetric group $S_3$ and $N \cong \mathbb{Z}/3$ is the normal
subgroup generated by a $3$-cycle. Consider $\mathbb{Z}$ equipped with
the trivial $G$-module structure. Determine the $G/N$-module structure
of the
\[
H^k(N,\mathbb{Z}).
\]
\end{dang}
Solution:


